
\subsection*{5.1}
\phantomsection
\label{sga3:VIIB:5.1}

Let \(k\) be a pseudocompact ring.\footnote{%
\label{sga3:VIIB:nde:143}The original assumes in \SGARefLink{sga3:VIIB:5.1:proposition}{5.1} that \(k\)
is a field. In fact, that hypothesis can be replaced by flatness
hypotheses; the editors modified \SGARefLink{sga3:VIIB:5.1}{5.1}--\SGARefLink{sga3:VIIB:5.1.5}{5.1.5} accordingly.}
Let \(G\) be a topologically flat infinitesimal formal
\(k\)-group, let \(A\) be its affine algebra, and let
\(B\subset A\) be a closed subalgebra. Put
\[
  X=\operatorname{Spf}(B),
\]
and let
\[
  r:G\longrightarrow X
\]
be the epimorphism induced by the inclusion \(B\subset A\). We
shall determine when \(r\) makes \(X\) the quotient of \(G\) on
the right by a subgroup \(H\)
\textup{(cf.\ \SGARefLink{sga3:VIIB:2.4}{2.4})}.\footnote{%
\label{sga3:VIIB:nde:144}The editors replaced ``left'' by
``right'' and modified the statement of \SGARefLink{sga3:VIIB:5.1:proposition}{Proposition~5.1} so as to
display both the equivalent conditions \textup{\SGARefLink{sga3:vii:paragraph:5-1:equation:i}{(i)}}, \textup{\SGARefLink{sga3:vii:paragraph:5-1:equation:ii}{(ii)}}
and the conclusion \(\operatorname{Spf}(B)\simeq G/H\) more
clearly.}

\medskip
\noindent\textbf{Proposition 5.1.}
\phantomsection
\label{sga3:VIIB:5.1:proposition}
{\itshape Let \(G\) be a topologically flat infinitesimal formal
\(k\)-group, let \(A\) be its affine algebra, let \(I_A\) be the
augmentation ideal of \(A\), and let \(B\) be a closed subalgebra
of \(A\). Put
\[
  I_B=B\cap I_A,
  \qquad
  J_B=AI_B.
\]
Assume that \(A\), and also
\[
  J_B^n/J_B^{n+1}
\]
for every \(n\geq0\), are topologically flat over \(k\). Then the
following two assertions are equivalent:
\begin{enumerate}
\item[\textup{(i)}] For every \(x\in B\),
\[
  \Delta_A(x)-x\widehat\otimes1
  \in
  A\widehat\otimes_kJ_B.
\]

\item[\textup{\SGAInlineRefTarget{sga3:VIIB:5.1:proposition:item:ii}(ii)}] Let
\[
  \pi:A\longrightarrow A/J_B
\]
be the projection, and put
\[
  \tau_1(a)=a\widehat\otimes1.
\]
The sequence
\[
  B\longrightarrow A
  \mathrel{\mathop{\rightrightarrows}^{\tau_1}_{%
    (\operatorname{id}_A\widehat\otimes\pi)\Delta_A}}
  A\widehat\otimes_k(A/J_B)
  \tag{*}\label{sga3:vii:proposition:5-1:equation:star}
\]
is exact. Equivalently, \(B\) is the set of all \(x\in A\) for
which
\[
  \Delta_A(x)-x\widehat\otimes1
  \in
  A\widehat\otimes_kJ_B.
\]
\end{enumerate}
In this case,
\[
  H=\operatorname{Spf}(A/J_B)
\]
is a formal subgroup of \(G\), and the sequence
\[
  G\times H
  \mathrel{\mathop{\rightrightarrows}^{\operatorname{pr}_1}_{\lambda}}
  G
  \longrightarrow
  \operatorname{Spf}(B)
  \tag{**}\label{sga3:vii:proposition:5-1:equation:double-star}
\]
is exact, where \(\lambda\) is the restriction of the
multiplication of \(G\) to \(G\times H\). Thus
\[
  \operatorname{Spf}(B)\simeq G/H.
\]}

\phantomsection
\label{sga3:VIIB:5.1:source-arrow-note}
\noindent\textit{Source note.}\enspace
In \textup{(*)}, the French PDF labels the lower arrow
\((\pi\widehat\otimes\operatorname{id}_A)\Delta_A\). The displayed
equivalent condition immediately before and after the sequence,
as well as the stated right-quotient convention, require
\((\operatorname{id}_A\widehat\otimes\pi)\Delta_A\), which is used
here.

\medskip
\noindent\textit{Proof.}
Put
\[
  \overline A=A/J_B,
  \qquad
  H=\operatorname{Spf}(\overline A).
\]
Then \(H\) is a formal subvariety of \(G\). Since
\(J_B\subset I_A\), the augmentation \(\varepsilon_A\) induces a
continuous \(k\)-algebra morphism
\[
  \overline\varepsilon:\overline A\longrightarrow k.
\]

If \textup{\SGARefLink{sga3:vii:paragraph:5-1:equation:i}{(i)}} holds, then
\[
  \Delta_A(I_B)
  \subset
  I_B\widehat\otimes1
  +
  A\widehat\otimes J_B.
\]
Therefore \(\Delta_A\) induces, after passage to the quotient, a
diagonal morphism
\[
  \overline\Delta:\overline A
  \longrightarrow
  \overline A\widehat\otimes_k\overline A.
\]
The maps \(\overline\Delta\) and \(\overline\varepsilon\) make
\(H\) a formal submonoid of \(G\). Since \(G\) is infinitesimal,
so is \(H\); \SGARefLink{sga3:VIIB:2.7:prop}{Proposition~2.7} then shows that \(H\) is a formal
subgroup of \(G\). By the definition of \(G/H\)
\textup{(cf.\ \SGARefLink{sga3:VIIB:2.4}{2.4})}, assertion \textup{\SGARefLink{sga3:vii:paragraph:5-1:equation:ii}{(ii)}} implies the final
assertion of the proposition.

It is clear that \textup{\SGARefLink{sga3:vii:paragraph:5-1:equation:ii}{(ii)}} implies \textup{\SGARefLink{sga3:vii:paragraph:5-1:equation:i}{(i)}}. The proof of
the converse occupies \SGARefLink{sga3:VIIB:5.1.1}{5.1.1}--\SGARefLink{sga3:VIIB:5.1.5}{5.1.5}.
